\chapter{Rigid Motion: Preserving the Invariant}

\otbookepigraph{Rotations are just a way of looking at the same world from a different angle without changing its distances.}{OT Interpretation}

\begin{chapteroverview}
    We de-construct the rotation matrix from a box of numbers into a structural guardian of geometry.
\end{chapteroverview}

\section{The Formal De-construction}
A rotation matrix \sym{rot-matrix} in \otindex{SO(3)} is defined such that:
\begin{equation}
    \sym{rot-matrix}^T \sym{rot-matrix} = I, \quad \det(\sym{rot-matrix}) = 1
\end{equation}

\section{The Re-interpretation: The ``Ruler'' Analogy}
\begin{intuitionbox}[The Rigid Ruler]
    Think of a rigid ruler. If you rotate it, the length marked on it does not change. Mathematically, the inner product \(\mathbf{x}^T \mathbf{y}\) must be invariant.
    \begin{equation}
        (\sym{rot-matrix}\mathbf{x})^T (\sym{rot-matrix}\mathbf{y}) = \mathbf{x}^T (\sym{rot-matrix}^T \sym{rot-matrix}) \mathbf{y} = \mathbf{x}^T \mathbf{y}
    \end{equation}
    This is \textit{why} the transpose must be the inverse. The geometry \textit{forces} the matrix to have this property.
\end{intuitionbox}

\section{Mental Re-construction}
\begin{practice}[Derivation Anchor]
    Close the book. Draw two frames \(o_0\) and \(o_1\). Write the coordinates of a unit vector \(\hat{i}_1\) as seen from frame 0. Why are the columns of \sym{rot-matrix} actually the unit axes of the new frame?
\end{practice}

\chapter{Forward Kinematics: The Geometric Chain}

\section{De-construction: The DH Convention}
The \acr{dh} convention uses four parameters (\(a, \alpha, d, \theta\)) to describe the transformation between links.

\section{Re-interpretation: The ``Joint-to-Joint'' Handshake}
\begin{designrule}[Chain of Command]
    Don't see forward kinematics as a single massive matrix. See it as a series of handshakes. Frame \(i\) says to Frame \(i+1\): ``Move along my common normal, then rotate about your own axis.''
\end{designrule}

\section{Stress-test}
\begin{reflectionbox}
    What happens if two joint axes are parallel? Does the \acr{dh} convention's unique common normal still exist? This is where the model is most likely to confuse the beginner.
\end{reflectionbox}
